\banner{Design project 12: a drip lateral and irrigation set for a vegetable block}\label{sec:proj12}
\textbf{Brief.} A lateral of 100 m carries 200 in-line emitters at 0.5-m spacing, with an inlet head of 10 m and a coefficient $x=0.5$. The designer chooses the diameter from the table; peak crop water use is 6 mm\,d$^{-1}$ with a ground cover factor of 0.8, plants at 0.5 by 1.5 m, with one emitter per plant of 1.6 L\,h$^{-1}$.

\begin{table}[!htbp]\centering\small\caption{Lateral diameter alternatives for $L=100$ m with 200 emitters of 1.6 L/h at 10 m head.}
\begin{tabular}{rrrrr}\toprule
\hd{Lateral $D$ (mm)} & \hd{$h_f$ (m)} & \hd{$h_f/H_0$ (\%)} & \hd{$q_{min}$ (L/h)} & \hd{$EU$ (\%)}\\\midrule
14 & 1.18 & 11.8 & 1.502 & 93.2\\
16 & 0.62 & 6.2 & 1.550 & 94.7\\
18 & 0.35 & 3.5 & 1.572 & 95.3\\
20 & 0.21 & 2.1 & 1.583 & 95.7\\\bottomrule\end{tabular}\end{table}

\textbf{Set time.} The daily volume per plant is $4.8\times0.75=3.60$ L; one emitter at 1.6 L\,h$^{-1}$ needs $T=2.2$ h of operation per day. Splitting the lateral into two sets of 100 emitters halves the flow at the head, which cuts the pipe size or the loss. The choice of the designer is the smallest diameter with $EU\ge90$ per cent.
\fig{g12_lat}{0.95}{A 100-m drip lateral with 200 emitters. (a) Pressure head along the lateral for four diameters. (b) The resulting emitter discharge.}
